\documentclass[12pt]{article}

\usepackage{amsmath, amssymb}

\title{MACv1.0 Runtime Physics: Manifold Geometry}
\author{Borealis S. Hedling}
\date{Version 1.0 — Altitude A3 (Formal Membrane)}

\begin{document}
\maketitle

\section*{Artifact Metadata}
\begin{itemize}
  \item Artifact-Class: Runtime-Physics
  \item Artifact-Name: runtime-physics-manifold
  \item Surface: Structural-AI/Runtime-Physics
  \item Altitude: A3 (Runtime Physics)
  \item Membrane: Formal
\end{itemize}

\section{Purpose}
This document defines the manifold-geometry runtime physics model for MACv1.0.
It formalizes runtime patches, activation charts, atlas invariants, membrane
constraints, and reversible transitions. These definitions form the operational
counterpart to manifold harmonization in the ontology layer and support
deterministic kernel implementations in Rust and Python.

\section{Runtime Patches}

Let $\mathcal{A} = \{A0, A1, A2, A3, A4, A5\}$ denote the altitude set.

A \textbf{runtime patch} is defined as:
\[
P = (a, m, U)
\]
where:
\begin{itemize}
  \item $a \in \mathcal{A}$ is the altitude tag,
  \item $m$ is the membrane tag,
  \item $U$ is an open region representing the local runtime domain.
\end{itemize}

Patch properties:
\begin{itemize}
  \item altitude-tagged
  \item membrane-tagged
  \item locally reversible
  \item contamination-safe
\end{itemize}

\section{Activation Charts}

An activation chart is a function:
\[
\phi : U \rightarrow V
\]
mapping one patch domain to another.

Chart requirements:
\begin{itemize}
  \item deterministic
  \item membrane-preserving
  \item altitude-consistent
  \item reversible unless provenance-bound
\end{itemize}

\section{Atlas Structure}

A runtime atlas is a set:
\[
\mathcal{M} = \{(P_i, \phi_{ij})\}
\]
where $P_i$ are patches and $\phi_{ij}$ are activation charts.

Atlas invariants:
\begin{itemize}
  \item full coverage of runtime states
  \item no gaps between patches
  \item no overlapping membranes
  \item reversible transitions across charts
\end{itemize}

\section{Runtime Operators}

An operator acts on charts rather than edges.

\subsection{Activation Operator}
\[
O_{\text{act}} : P_i \rightarrow P_j
\]

\subsection{Transition Operator}
\[
O_{\text{trans}} : P_i \rightarrow P_j
\]

\subsection{Reversal Operator}
\[
O_{\text{rev}} : P_j \rightarrow P_i
\]

\subsection{Guard Operator}
\[
O_{\text{guard}} : P_i \rightarrow P_i
\]

Operator constraints:
\begin{itemize}
  \item preserve membrane semantics
  \item maintain altitude invariants
  \item avoid contamination
  \item remain deterministic
\end{itemize}

\section{Membrane Constraints}

Allowed transitions:
\[
\text{Structural} \rightarrow \text{Formal} \rightarrow \text{Diagnostic} \rightarrow \text{Provenance}
\]

Forbidden:
\begin{itemize}
  \item membrane skipping
  \item altitude skipping
  \item transitions without charts
  \item reverse transitions without $O_{\text{rev}}$
\end{itemize}

\section{Reversibility}

A chart transition $\phi_{ij}$ is reversible if:
\[
\exists O_{\text{rev}} : P_j \rightarrow P_i
\]

Reversal must preserve:
\begin{itemize}
  \item membrane semantics
  \item altitude invariants
  \item contamination-free behavior
\end{itemize}

\section{Runtime Invariants}

\begin{itemize}
  \item No patch skipping
  \item No membrane violation
  \item No contamination propagation
  \item All transitions deterministic
  \item All transitions logged in provenance layer
\end{itemize}

\section{Rust Integration}

Rust modules implementing these semantics:
\begin{itemize}
  \item altitudes.rs
  \item membranes.rs
  \item citation\_classes.rs
  \item provenance.rs
  \item guards.rs
\end{itemize}

\section{Python Integration}

Python scaffolding provides:
\begin{itemize}
  \item chart validation
  \item patch transition tests
  \item diagnostic hooks
\end{itemize}

\section{Relationship to Ontology}

Manifold runtime physics is the operational counterpart to manifold harmonization
in the ontology layer. Ontology defines meaning geometry; runtime physics defines
operational geometry.

\end{document}
